\section{Is the Black Mass Valid}

\label{sec:BlackMass}
by \textbf{Solange Herz}

A frustrated wag in these confused times is said to have asked his Bishop, ``Your Excellency, does attendance at the black Mass satisfy my Sunday obligation?''

In view of the ceremonies now permitted in many Catholic churches, the question may not be as flippant as it sounds. It may come closer to penetrating the workings of iniquity than many a learned treatise.

Given the current aberrations, the query is sound, throwing the modern apostasy into disturbing perspective. The Gospels tell us that when our Lord was casting out devils, they testified in their rage, ``Thou art the Son of God'' (Luke 4: 41), much as the corrupt prophet Balaam foretold the Incarnation while attempting to curse Israel. In the same way, the black Mass gives witness to the true Sacrifice as surely as the Father of Lies himself told the truth when he promised Eve that she and Adam would not die but in due time be as gods. It so clamors for attention today and a search into its dark origins may tell us something helpful.

St. Justin Martyr, a convert from paganism, wrote a famous description for the Roman Emperors Marcus Aurelius and Lucius Verus, of the Christian Mass as it was celebrated in the second century. He explains therein, ``We call this food `eucharist,' and none may partake of it if he believes not in the truth of our doctrine, has not received the washing for the forgiveness of sins and regeneration, and lives not according to the precepts of Christ. For we do not take this food as common bread and common drink, but just as our Savior Jesus Christ, incarnate by the power of the Word of God, took flesh and blood for our salvation, so the food consecrated by the prayer formed of the words of Christ, this food which is to nourish our blood and flesh by being assimilated, is the flesh and blood of the incarnate Jesus: such is our teaching. For the Apostles, in their memoirs called Gospels, report that Jesus gave them these instructions: ``Do this in memory of me, this is My Body''; He likewise took the cup, and having given thanks He said: ``This is My Blood'': and He gave it to them only.'' Whereupon St. Justin adds, ``This the evil spirits have imitated by instituting the mysteries of Mithra; for you know, or may know, bread and a cup of water are given in the ceremonies of initiation, and certain formulas are pronounced.'' It would be forcing the text to assert that the black Mass as we know it today is referred to here, but it is here in essence, on its way to further development as Satan's power over men increased.

The Mass of Mithra was already a product of long preparation, for the satanic Mass actually began when Adam and Eve partook of the forbidden fruit. This was, historically, the First Communion of persons hoping to become gods on their own with only the devil's help, and it set the pattern for all that followed.

It was soon ritualized by the murderous Cain, figure of the reprobate, who offered ``of the fruits of the earth'' gifts for which the Lord ``had no respect'' (Gen. 4:3, 5). It was therefore a parody of true worship, springing not from heaven, but from the earth and man's own vitiated will. The gifts Cain offered were undoubtedly good in themselves, but the hidden intention of the sacrifice, known only to God, rendered them unacceptable, ``valid'' as they were under the first natural dispensation.

We have reason to believe that Cain's liturgy proliferated, however, for Noah's sacrifice after the Flood was simply the reestablishment of the slain Abel's approved worship, destroyed from the face of the earth by the sons of Cain, who had thereby brought the disaster upon the whole world. Noah's was adoration according to righted natural law, followed in due course by Abraham's divinely inspired offering of his son Isaac, the promise of the real Sacrifice to come, ritually prefigured by Moses and instituted in reality by our Lord before His Passion. Outside this line of development, no worship acceptable to God can exist. He Himself had to come and show us how He wished to be worshipped, giving Himself to us as both Priest and Victim, no others being proportionate to such an undertaking.

There is no such thing as purely human religion. As the Psalmist said, ``All the gods of the gentiles are devils'' (Ps. 95:5). ``The things which the heathens sacrifice,'' St. Paul tell us, ``they sacrifice to devils, and not to God,'' regardless of what the heathen may profess or believe. ``And I would not that you should be made partakers with devils. You cannot drink the chalice of the Lord, and the chalice of devils: you cannot be partakers of the table of the Lord, and of the table of devils'' through some misguided ecumenism (1 Cor. 10: 20-2l).

Eusebius and other Fathers of the Church taught that demons instructed men in the black arts, because it seems evident that men could never have discovered these by themselves, any more than they could have contrived the genuine rites of divine institution. Devils have always pretended to be gods in order to be worshipped. If they have one consuming interest, it is liturgy. Behind every form of true worship has lurked the somber shadow of the black Mass, keeping pace with the mysteries of Redemption, even as the envious Cain skulked behind Abel, luring him onto open ground in order to kill him.

Coexistent with the Church and necessarily developing in her wake, the black Mass feeds on faith, for no unbeliever would think of defying or bypassing God, let alone adoring His primordial Enemy. Hard-core Satanists deplore the loss of ``faith'' among colleagues who regress into worship of mere Nature or other substitutes. Historically, however, it can be noted that occult practices rise and fall with rationalism.

This may seem a contradiction, but not when we understand that witchcraft and rationalism are manifestations of the same disease: naturalism, applied in one case to religion and in the other to intellectual life. The devil, although a spiritual being, is so by nature, and he is bound to his purely natural existence.

He seeks to destroy the supernatural life of grace in man by deflecting supernatural worship of God first into mere natural worship of God and thence into worship of man, culminating in the worship of his satanic self. Within this progression lie all shades of what we euphemistically label heresy.

The whole body of witchcraft is a monument of perverted rubrics. Like Pharaoh's magicians who imitated all Moses' signs and wonders, devils have produced spells and incantations in lieu of blessings and prayers for every occasion, amulets and charms for sacramentals. For true mysticism they propose delusions, euphoric and otherwise, most evident in our day in Pentecostalism and various ``eastern'' disciplines. Sacraments of initiation aping Baptism abound. Even Chinese Communists have one.

Everywhere natural is substituted for supernatural, and carnal communication of some kind, often in spiritual guise, replaces union with God in spirit and truth.

Dismissing the Church's vast literature on witchcraft as superstitious legend piously believed by ingenuous churchmen in medieval times is itself a bewitchment of the devil, who would readily have us think he doesn't exist, until we have become sufficiently weakened. Not only profoundly unscientific in the face of today's mounting evidence, skepticism is blasphemous, for God Himself attests to the reality of sorcery:

``Wizards thou shall not suffer to live'' (Ex. 22: l8). ``A man or woman, in whom there is a pythonical or divining spirit, dying let them die'' (Lev. 20:17). ``The soul that shall go aside after magicians and soothsayers, and shall commit fornication with them, I will set my face against that soul'' (Lev. 20:6). Charmers, fortune-tellers, and those ``seeking truth from the dead'' are likewise anathematized, as are those indulging in ritual sacrifice of children (Deut. l8: 10-11). Like the Old Testament prophets before them, St. Peter and St. Paul often found themselves pitted against occultists, as the Acts of the Apostles testify.

The Church soon translated this into canon law, the first juridical mention of witchcraft extant being a decision of the Council of Ancyra in 314 recorded in Gratian's Decretum. St. Thomas viewed sorcery seriously, even recognizing that it could set up a valid impediment to the Sacrament of Matrimony by inhibiting the power of sexual intercourse. No reputable student of occultism could overlook the Malleus Maleficiarum, the ecclesiastical handbook on witchcraft put together in 1486 for the guidance of Inquisitors, together with the wealth of other works which it inspired.

In a letter a century before, Pope John XXII had expressed fears for his life at the hands of witches, and during the l4th and 15th centuries almost every Pope issued bulls against witchcraft. These culminated in the Great Bull Sumnis Desiderantes of Innocent VIII in 1484, issued at the insistence of the two Dominicans, Jacob Sprenger and Heinrich Kramer, who compiled the Malleus two years later. Listing in some detail the heinous practices of witches, this authoritative document notes, ``Above and beyond this, they blasphemously renounce that Faith which they received by the Sacrament of Baptism, and at the instigation of the Enemy of the human race they do not shrink from committing and perpetrating the foulest abominations and excesses to the peril of their souls, whereby they offend the Divine Majesty and are a cause of scandal and dangerous example to very many.'' Let him dismiss this as pious nonsense who dares.

Sober study reveals that devils are highly organized. Like the Church it apes, the ``mystical body of Satan'' has its structures, specialists, ministers, apostles and theologians, as it has its faithful. At its center there must of course be a supreme liturgical act, as in God's Church. This can only be what we call the black Mass, a rite whereby the Mysterium fidei is formally supplanted by the myaterium iniquitatis, ritual ratification of its devotees' fall into the Great Temptation.

The black Mass is the supreme fruit of witchcraft, to which all its history and paraphernalia tend. As Holy Mass is ordered to union with Christ, the diabolic Mass is directed to a similar close, personal union with the devil. Its effects far worse than mere possession, which is ordinarily limited to the body, the black Mass engages the human will and all the spiritual faculties --- although necessarily only on the natural level, supernatural action being impossible without divine grace. Immolation, however, is essential, for this specifically binds the will.

\begin{center}* * *\end{center}

The Mass of Mithra that St. Justin speaks of celebrates a conflict between good and evil by which the divine huntsman Mithra, a personification of light and goodness, effects the salvation of the world by pursuing a bull personifying darkness and evil into a cave, where he dispatches him with a sword on orders from the Sun-god Ormuzd --- whereupon all visible material creation springs from the bull's dead body. (We might mention that Mithra wears a Phrygian cap like those sported by the French revolutionists, a cap which can also be seen on the heads of the figures of Liberty and Young America in the blasphemous ``Apotheosis of Washington'' painted by Brumidi in the dome of the U.S. Capitol.) Supported by the Emperors, Mithraic worship spread along the Euphrates and through Europe along the Rhine and Danube. Even after the public cult ceased with Constantine, it flourished in England, where the Roman legions had built temples from London to York, so that it became known as the Island of the Bull --- today, ``John Bull.'' Hugh Ross Williamson in The Arrow and the Sword even voices the suspicion that the butchery of St. Thomas Becket in his cathedral may have been a ritual murder in the Mithraic tradition. As in Masonry, only men were devotees, and their initiation took the form of a ``baptism'' simulating death and resurrection, more than vaguely reminiscent of the Masonic initiation ceremony of the ``raising of Hiram.''

Sacrifice of a bull constituted the main rite of worship, vestiges of which are still found in Germany today, and it's curious to note that after the Masonic ceremony dedicating the U.S. Capitol, the participants feasted on roast ox. Mithra's adepts likewise communicated in the common meal mentioned by St. Justin, commemorating the one he shared with the Sun-god after killing the bull. An ancient Mithraic hymn reads, ``Thou hast redeemed us too by shedding the eternal blood,'' for the task of man was to liberate his soul from the material shackles of his body and follow Mithra to a heaven of pure spirit. With its Christlike savior and its concept of sacrifice leading to resurrection, this myth came dangerously close to Christian doctrine and misled many. St. Augustine tells of a priest of Mithra he knew who used to say, ``Our capped one is himself a Christian.'' On the contrary, like all mystery cults, Mithraism was in no sense of Christian Inspiration, but pure Gnosticism. Derived originally from the Jewish kabbala imported from Egypt and Babylonia by heretical Jews, gnostic beliefs of one kind or another had been at work adulterating the true religion from earliest times, the prophets inveighing against them to little avail. Systematized 200 years before Christ in a blend of Hebraic and heraetic culture by Hellenistic Jews like Philo of Alexandria, Gnosticism by the time of St. Justin was presuming to replace not only polytheism and Judaism, but also Christianity, by allegedly combining in itself the higher principles of all three. It was, in other words, an early attempt at ``world religion.''

As its name implies, it proposes salvation by imparting superior knowledge or gnosis, the primordial ``ancient wisdom'' revealed secretly to a chosen few from generation to generation. Vaguely pantheistic, it is necessarily allied to some form of dualism, for Satan, who inspired it, considers himself equal to God. In order to present himself as a viable candidate for the Lordship of the Universe, he must establish himself on the same divine level as God. He sought to do this at the very beginning of human history in Eden, where -- as Louis Veuillot so aptly put it -- he sought to demote the sovereign God to the rank of simple citizen in paradise according to democratic principles. At the same time, he imparted the basic tenet of the gnosis, a special knowledge not given by God to man: He asks Eve, ``Why hath God commanded you, that you should not eat of every tree of paradise?'' When Eve replies that to eat the fruit of a certain one entails the death penalty, the serpent tells her, ``No, you shall not die the death. For God doth know that in what day soever you shall eat thereof, your eyes shall be opened and you shall be as Gods, knowing good and evil'' (Gen. 3:1-5).

She would be ``enlightened.'' It is precisely this approach to good and evil which constitutes Gnosticism. Whereas Christian faith tells us there is only one God, Who is all good, eternal and almighty, gnostic dualism would have us believe the world is governed by two equal but warring principles, one good and the other evil, each limiting the power of the other: God and Satan. The spiritual world allegedly sprang from the good principle, but the material world, which sprang from the evil principle, as from Mithra's bull, is intrinsically evil.

To say that evil comes from on high and has an existence of its own, rather than being the effect of personal sin, is a denial of free will and makes pessimism an article of faith. We can see how deeply gnosticism had infected Adam from his ``communion'' with Eve, when later he has the temerity to tell God that he sinned because of ``the woman, whom Thou gavest me to be my companion'' (Gen. 3: 12). Like all gnostics, Adam was ascribing evil to God, who made him sin. Christianity of course teaches that man's will is free, and that everything without exception created by ``the Father almighty, Maker of heaven and earth and of all things visible and invisible'' is intrinsically good. And this includes matter. The Christian's help is in the name of the Lord Who created both heaven and earth. His formal Credo couples ``the resurrection of the body'' with ``life everlasting'' in the selfsame sentence, at the close of both the Apostles' and the Nicene Creed. And of course God also created Satan, who, although originally good, became corrupt only by his own will.

Gnosticism is not merely the primordial heresy, generating all others. It is the very taproot of Satanism, for by a total reversal of truth, it proceeds to teach that the good principle is not God, but Satan.

It is Satan who created the good, purely spiritual world to which man must aspire to be saved: whereas God, who passes for good, is actually the evil principle who created that evil, debasing substance called matter, from which men must free themselves. This they can do when they have acquired the special knowledge Satan gives them. This is what the serpent was insinuating to Adam and Eve. Worship of Satan is therefore not a new religion, says the gnostic, but the Old Religion, antedating the material world and driven underground by the usurper who calls Himself God.

According to this, the true goodness of the one Scripture calls the Serpent, the Tempter, the Adversary, is proved by the fact that he led the revolt against the Creator of the evil material world. The Serpent would therefore be the Logos, hero of the Revolution, making of Christ a devil who interferes with the true Jesus still to come. Unless carefully interpreted according to this secret wisdom, everything the Bible says is wrong, where characters like Cain, the rebels of Core and Judas Iscariot are portrayed as reprobates, when they are in fact calumniated saints. A second-century gnostic sect actually called themselves the Cainites and disseminated a ``gospel of Judas.''

It is terrifying to see how common these ancient ideas are becoming today in the current cult of the anti-hero in preparation for the coming of the AntiChrist. Judas now figures openly as the hero of popular musicals, and in ordinary subculture parlance the very word ``bad'' now means ``good.'' In such context, sin ``liberates'' and ``enlightens.'' It becomes a positive duty, a defiance of the ``evil'' God who has trapped us in His wicked material establishment. The black Mass becomes the supreme means of giving the devil his due, homage to the Bright Angel unjustly exiled from heaven, to whom great wrong was done.

``Woe to you that call evil good, and good evil: that put darkness for light, and light for darkness, that put bitter for sweet, and sweet for bitter!'' (Is. 5: 20).

It is no longer our Lord Jesus Christ, true God and true Man, God of God and Light of Light, Who came down from heaven to save us by dying on the Cross, but Lucifer, the Light-Bearer of darkness, who will be the light of the world. He burns to hear the eternal Sanctus sung before the throne of his own trinity, the dragon, the beast and the false prophet, his Benedictus deeming blessed he who comes, not in the name of the Lord, but in his own name. So did our Lord prophesy to the Jews: ``I am come in the name of my Father, and you receive me not: if another shall come in his own name, him you will receive'' (John 5: 43).

Satan, angelic spirit rotten with pride, cannot forgive God for having assumed human flesh from a woman and not becoming an angel like himself! It was to Hebrews tainted with gnosticism that St. Paul addressed the question, ``For to which of the angels hath (God) said at any time, `Thou art my son, today have I begotten thee'?'' (Heb.1:5). And it was into the very teeth of the mounting heresy that the Council of Nicea flung the words of its Creed: Et incarnatus est de Spiritu Sancto ex Maria virgine: ET HOMO FACTUS EST.

0 Mary, in destroying gnosticism, thou hast destroyed all heresies! As St. Jerome put it, ``The truth has set bounds.'' Being simple, it is quickly defined, ``but evil and falsehood multiply without end, and the more these are pursued, the more errors they produce.'' It would take several lifetimes to trace the labyrinthine developments of Gnosticism through the centuries to our day, let alone describe all its shapes. How can there be unity and simplicity in the works of one whose very name is Legion?

Eventually worming its way into the heart of Christianity by way of pagan metaphysics, Gnosticism has festered for centuries as a religion within a religion, all the while bearing the black Mass in its bosom.

From Mithraism it reappeared among the Novatians and thence to the Manichaeans, fueling so powerful a movement that even giant intellects, like St. Augustine, were ensnared in it for a time. It culminated soon in the major outbreak called Catharism, a rebellion against the Church which began in the eleventh century among some Bulgarian monks and spawned hundreds of sects variously known as Albigensians, Bogomils, Lombards, Valdenses, Mandeans, Johannites or whatever. Its fury and violence were such that God dispatched St. Dominic and his preaching friars everywhere to rescue Christendom, spiritually armed for the occasion by our Lady herself with the Rosary. The word of God and the Hail Mary setting the stage for action, the first Inquisition was authorized by the Holy See, with St. Dominic very probably its first Inquisitor. Nor was physical force neglected, for Simon de Montfort and his troops battled relentlessly in what they knew was a deadly struggle for thousands of imperiled souls.

Henry Taylor Fawkes Rhodes, an English criminologist who in 1968 published a study called The Satanic Mass, says the Cathars ``denied the validity of the Christian sacraments because they said that matter, being evil, could not be the vehicle of Divine Grace. One of their catechisms is said to have included the following questions and answers: `Do you renounce the cross which the priest made on you in your baptism on breast and shoulders and head with oil and chrism?' `I do renounce it.' `Do you believe that water works salvation for you?' `I do not believe it.' They of course rejected Communion and the Mass for the same reason.''

The reader is free to note what parallels may exist here with the changes now made in the administration of Catholic sacraments. ``On the other hand, there was a Cathar substitute for baptism where the officiant breathed seven times in the face of the candidate, saying, `Receive the ``Holy Spirit'' from ``Good Men.''\,' Their ``Communion'' was a commemorative meal of bread, fish and wine. ``Likewise the Cathars denounced the visible and material foundations of the Catholic Church as an evil creation and the Synagogue of Satan. They also rejected the cult of the Virgin Mary because they said that the Son of God could not have had so close an association with the material world as to be born of a woman.'' They discouraged marriage, high initiates practicing ``severe austerities, among which almost total abstinence from food and savage flagellation were popular.'' Highly organized, they had their own form of holy orders, but only the perfect were initiated into its deepest secrets. As the old occult adage has it, ``Nothing is concealed -- from him who knows!''

St. Bernard called Catharism ``the deceit of devils.'' According to Rhodes, ``To spread their doctrines among the orthodox, the Cathars used the method of infiltration. Outwardly they conformed to the Christian way of life and to Christian practices. They used the same formulae, but interpreted them in the opposite, Cathar, sense. It is this fact which substantiates the traditional belief that the Black Mass is an affair of `contrary rites and ceremonies' and the saying of prayers `backwards.' This cannot be taken in the literal sense, but essentially it describes the ambiguous position of men who said one thing and meant another... . It was possible for a priest to introduce prayers and ceremonies into an orthodox rite without any of the congregation, except the `Satanically' initiated, being aware that anything unusual was taking place.''

Catharism is generally credited with destroying the Knights Templar, an order whose members were under perpetual religious vows for the military defense of Christendom. Many of them came to wear the cordula, a little cord which formed part of the secret clothing of the sect. In the judicial inquiries, they were accused of defiling the Cross, denying the divinity of Christ and permitting laymen to give absolution, along with practicing sodomy and onanism in connection with their rites. Their ``Mass'' was devoid of the words of Consecration and therefore of sacramental Communion, but human sacrifice was often made use of instead. In 1312, although some members were still not infected, the Order had to be dissolved by Apostolic Decree. The evil persisted, however, and even a century later the Marshal of France, Gilles Rais, was convicted of murdering 140 children for such practices, some authorities placing the number closer to 800.

\begin{center}* * *\end{center}

A creation of the satanic intelligence, Gnosticism's theology exudes a power to be reckoned with. Clothing itself as it does in any religious symbolism which suits its purpose, and with no clear-cut depositum fidei to defend, it pretends to operate in upper regions where mere creeds are of no special importance. Its words can mean anything. Always rooted in pessimism, it offers escape from an intrinsically evil material world through mystic knowledge which antedates and engulfs all religious differences. It was plainly visible in the teachings of John Calvin, erupting in the Catholic Church as Jansenism. With Schopenhauer, it emerged as modem philosophy, where evil is presented as a positive principle, and good as simply a negation of evil, equated with ``well-being.'' Teilhardism, in which the spiritual principle of the world evolves and gradually liberates itself from matter into the full consciousness of a ``noosphere,'' is sheer Gnosticism, as was Hegel's philosophy before it. So is Communism, which sees a ``good'' Proletariat working its way out of its material bourgeois shackles into perfect spiritual freedom.

These errors are all found in Modernism, a plague to which St. Pius X in the encyclical Pascendi gave a definition which could be applied equally to Gnosticism: ``the synthesis of all heresies.'' He ably described its equivocal methods, so like Cathar ambiguity. So successful have they been that in 1968 an Archbishop Lefebvre could speak of the ``paralysis'' of the Church's magisterium, ``manifested by the lack of definition of notions and terminology, by the lack of preciseness, necessary distinctions, so that one no longer knows what it means to speak: think of those words `human dignity,' `liberty,' `social justice,' `peace,' `conscience'... . From now on, in the Church itself, these words may be given either a Marxist or a Christian meaning with equal conviction.''

\begin{center}* * *\end{center}

How did this happen? The phenomenal acceleration of Gnosticism since the ``enlightenment'' of the eighteenth century can only be explained by the emergence of speculative Masonry, soon aided by a whole network of secret gnostic societies, moving, multiplying, dividing, colonizing and passing in and out of sight like microbes under a microscope --- groups like the Golden Dawn, the Ordo Templi Orientis, Stella Matutina, the Hermetic Brotherhood of Luxor, Theosophy, Anthroposophy, etc.

In Israel et L'Humanite, Rabbi Benamozegh says, ``What is certain is that Masonic theology corresponds well enough to that of the kabbala,'' with the following editorial note added by his editors, Dr. Modiano and the Grand Rabbi Toaffi: ``To those who may be surprised by the use of such an expression, we would say that there is a Masonic theology in the sense that there exists in Freemasonry a secret, philosophic and religious doctrine, which was introduced by the Gnostic Rosicrucians at the time of their union with the Free Masons in 1717. This secret doctrine, or gnosis, belongs exclusively to the High, or philosophic, degrees of Freemasonry.''

``Among all the ancient nations there was one faith and one idea of Deity for the enlightened and educated, and another for the common people,'' wrote Albert Pike, high priest of Masonry in Morals and Dogma. ``To this rule the Hebrews were no exception. Lady Queenborough, in Occult Theocrasy, calls Gnosticism the ``mother'' of Freemasonry, which ``imposed its mark in the very center of the chief symbol of this association.'' ``The most conspicuous emblem which one notices on entering a Masonic temple, the one which figures on the seals, on the rituals, everywhere in fact, appears in the middle of the interlaced square and compass, is the five-pointed star framing the letter G... . In the lower grades, one is taught that it signifies Geometry. To the brothers frequenting the lodges admitting women as members, it is revealed that the mystic letter means Generation, but the revelation is attended with great secrecy. Finally, to those found worthy to penetrate into the sanctuary of Knights Kadosch, the enigmatic letter becomes the initial of the doctrine of the perfect initiates which is Gnosticism... defined by Albert Pike as `the soul and marrow of Freemasonry.'\,''

In 1925 one of Masonry's most accomplished mystical theologians, Jirah Dewey Buck, set out to establish this in an authoritative work entitled Mystic Masonry and the Greater Mysteries of Antiquity. He speaks of ``the cardinal tenets of the old primitive faith, which underlie and are the foundation of all religion,'' maintaining, ``There is no fact in history more easily and completely demonstrable than the existence of the Secret Doctrine in all ages among all people, and of Adepts or Masters who were familiar with its teachings, and were more or less capable of expounding its mysteries. It is equally demonstrable that this Secret Doctrine was the real foundation of every great Religion known to man; that only the initiated Priest or Hierophant knew the real doctrines in any case... .''

Furthermore, ``the sacred books of all religions, including those of the Jews and the Christians, were and are no more than parables and allegories of the real Secret Doctrine, transcribed for the ignorant and superstitious masses.'' Underlying this Secret Doctrine, ``which came from the East,'' was a ``profound philosophy of the creation or evolution of worlds and of man. The present humanity, in many quarters of the globe, has evolved on the intellectual plane so far that there now exists a very large number of persons capable of apprehending this old philosophy, and at the same time, capable of understanding the responsibility incurred in misusing or misinterpreting it. A large number of persons have reached, on the intellectual plane, the state of manhood and are capable of partaking of the `fruit of the tree of knowledge of Good and Evil'... . It is therefore high time that the philosophy of the East should illumine the science of the West, and thus give the death blow to that intellectual diabolism and spiritual nihilism known as Materialism, and this only the Secret Doctrine can accomplish.'' No Cathar could have put it better.

Elsewhere Buck maintains that ``for centuries designing Priests, many of whom would have disgraced a scaffold, but who have been canonized as saints, have done their utmost to deprive humanity of this knowledge... . There is a Grand Science known as Magic, and every real Master is a Magician. Masonry in its deeper meaning and recondite Mysteries constitutes and possesses this Science, and all genuine Initiation consists in an orderly unfolding of the natural powers of the neophyte, so that he shall become the very thing he desires to possess.''

Additionally, ``Drawn from the Kabalah, and taking the Jewish or Christian verbiage or symbols, Masonry discerns in them universal truths, which it recognizes in all other religions. Many degrees have been Christianized only to perish... circumscribed by narrow creeds... . Masonry is the Universal Religion only because and only so long as it embraces all religions... . It is universal and eternal... . No anathemas of Popes so lasting as to count one second on its Dial of Time! These, one and all, serve only to keep the people in darkness and retard the reign of Universal Brotherhood.''

Brother Buck tells us that proofs of the diffusion of the Secret Doctrine, ``authentic records of its history, a complete chain of documents, showing its character and presence in every land, together with the teaching of the great Adepts, exist to this day in the secret crypts of libraries belonging to the Occult Fraternity.'' He says the ``Great Republic, composed of many Nations and all people,'' was ``a cardinal doctrine in the Ancient Mysteries'' and ``the logical deduction from our idea of Divinity, and of the essential nature and meaning of Christos. Humanity, in toto,... is the personification of the Divine in Creation.'' In other words, man is God.

He concludes: ``Drop the barnacles from the Religion of Jesus, as taught by Him, and by the Essenes and Gnostics of the first centuries, and it becomes Masonry. Masonry in its purity, derived as it is from the old Hebrew Kabalah, as part of the Great Universal Wisdom-Religion of remotest Antiquity, stands squarely for the Unqualified and Universal Brotherhood of Man, in all time and in every age.'' Such words bring into sharp relief Leo XIII's encyclical Humanum Genus, against the Masonic Brotherhood, whose ``ultimate aim is to uproot completely the whole religious and political order of the world which has been brought into existence by Christianity and to replace it by another in harmony with their way of thinking.'' The Holy Father speaks advisedly of the danger to the political order, as well as to the religious, for Gnosticism has always plotted the destruction of both. As we shall see, the history of the black Mass, not usually envisaged as a political tool, would nevertheless prove it.

Henry Rhodes writes: ``In regarding the material universe as a creation of evil, the Cathars included human society in the same condemnation in which they involved the Christian Church. Their aim was therefore to overthrow the existing order and to substitute a new one on the Cathar model. The severe persecution they suffered derived from the belief that they were a danger not only to religion, but to the foundations of society itself. It was a necessary part of their theory to denounce the pomps and vanities of the world, and the riches and power of priest and noble in much the same terms that the modem revolutionary denounces economic powers vested in the hands of the few.''

For the gnostic, Satan is something more than a god. Driven from heaven and forced to lie low in dark corners of Christendom, he becomes rather the Great Underdog with whom dissatisfied mortals laboring under real or imagined oppression can readily identify. Polarizing malcontents, he becomes their champion against God Himself, Who is made responsible for their sorry lot first by creating their cruel material world, and then by refusing them equality, freedom of conscience and unlimited wealth. Some theorists, like the liberal French historian Jules Michelet, have maintained that Satanists and sorcerers in any age were never anything but the anarchists and nihilists of their day. True it is that the occult arts have always flourished in times of political upheaval, for witchcraft is deeply social in character, the very kettle of Satan's ongoing Revolution. The primitive black Mass, first properly discerned among the Cathars, developed elaborate ritual in the Reformation, rising to greater heights at the time of the French Revolution, until our day, when churches for its celebration are openly erected on the streets of Megalopolis.

\paragraph{Conclusion} IN 1952 THE ENGLISH CLERGYMAN Walton Hannah in DARKNESS VISIBLE published documentation of an elaborate ritual called the ``Ceremony of the Rose-Croix of Heredom,'' printed privately in 1926 for the 33rd degree Supreme Council. At its conclusion a biscuit is given by the ``Most Worshipful Sovereign'' to the ``Prelate'' which is broken, dipped in salt and eaten by those present, followed in like manner by a chalice of wine. Whereupon the Sovereign says, ``All is consumed,'' to which the Prelate responds, ``Gloria in excelsis Deo, et in terra pax hominibus bonae voluntatis.'' Scholars find significant parallels between this service and the ancient Mass of the Vain Observance in use among the gnostic Knights Templar, where the words of consecration were purposely omitted. It is also similar, we might add, to the Consecration-less Mass concocted by Benjamin Franklin and his occultist friend Francis Dashwood, Lord le Despencer, erstwhile ``Prior'' of an unsavory community known as ``the mad monks of Medmenham'' which numbered the Earl of Sandwich and other libertines among its members. And there was the ``Adonaicide Mass'' composed not long after by the Jew Moses Holbrook, Grand Master at Charleston before Pike's ascendancy.

None of these liturgists, however, could compare with Aleister Crowley, the hard core Satanist popularly known as the Great Beast, who distinguished himself as a member of the Ordo Templi Orientalis (OTO) founded, in 1902, by Karl Kellner and inner circle Masonic adepts. Their purpose was to collaborate in developing ceremonials which would systematically communicate not only the secrets of Masonry proper, but also those of the Gnostic Catholic Church, the Rosicrucians and other satellites through a new system of ten degrees of ascending ritual. A talented writer whose prose and poetry found their way into respectable anthologies, Crowley composed an ode, in 1899, at the close of the Spanish-American War called ``An Appeal to the American Republic'' to celebrate the overthrow of ``priest-craft and tyranny.'' Needless to say, he set to work with a will perfecting the new ritual.

He tells us in his Confessions that he didn't always know all the variations of the secret grip, but as a friend of the Imperial Grand Hierophant, John Yarker, he became a kind of inspector general of Masonic rites ``charged with the secret mission of reporting on the possibility of reconstructing the entire edifice.''

So Masonry has its ``renewals'' too! He quotes a past Grand Master writing in the English Review for August, 1922, who deplored the general refusal of Freemasons, especially in England and America, to take their Masonry seriously.

This gentleman admits that the political orientation of lodges far outweighs the religious, Craft Lodges in England having originally figured as Hanoverian clubs, the Scottish as Jacobite, and Cagliosotro's Egyptian Lodges as revolutionary cells. But why then, this gentleman asks, if that's all there is to Masonry, whose rites are merely vestiges of initiations into old trade guilds, ``should such a matter be hedged about with so severe a wardenship, and why should the Central Sacrament partake of so awful and so unearthly a character?''.

Why indeed? He continues, ``As freemasonry has been `exposed' every few minutes for the last century or so, and as any layman can walk into a Masonic shop and buy the complete Rituals for a few pence, the only omissions being of no importance to our present point, it would be imbecile to pretend that the nature of the ceremonies of Craft masonry is in any sense a `mystery.' There is therefore no reason for refraining from the plain statement that, to anyone who understands the rudiments of symbolism, the Master's degree is identical with the Mass. This is in fact the real reason for the papal anathema: for freemasonry asserts that every man is himself the living, slain and re-risen Christ in his own person.

He adds, ``It is true that not one mason in ten thousand in England is aware of this fact; but he has only to remember his `raising' to realize the fundamental truth of the statement.'' It was for the OTO that Crowley eventually produced the ritual of the Gnostic Catholic Church, ``the central ceremony of its public and private celebration corresponding to the Mass of the Roman Catholic Church.'' Thus, at last, the ``black'' Mass is called by its true name: the Gnostic Mass. Anyone who can stomach pages of blasphemous parody of high literary quality may read the text as published in GEMS FROM THE EQUINOX, edited by Israel Regardie, in 1974. Its creed, recited by Deacon and People, begins, ``I believe in one secret and ineffable Lord: and in one Star in the Company of Stars of whose fire we are created, and to which we shall return.
... I believe in one Earth ... And I believe in the Serpent

and the Lion, Mystery of Mystery, in His name Baphomet. And I believe in one Gnostic and Catholic Church of Light, Life, Love and Liberty, etc.'' Needless to say Simon Magus is reverently mentioned in a Communicantes.

As expected in a celebration of dualism, both a priest and a priestess officiate. As Henry Rhodes noted, ``As representative of the material creative principle'' the substance of the sacrifice ``must be of male and female, and of sexual origin. However primitive and infantile this may appear, it is not obscenity for obscenity's sake,'' but the formalization of blasphemy against God's sole and unique Fatherhood as Creator.

This part of Crowley's Mass is set in mystical cipher, but he himself gives the clue when he says, ``Excrement is my host.'' When he died, in 1947, a judge characterized him as the most perverted man in all England, an allegation never challenged. Worse than that, however, he emerges as a first class liturgist, the veritable Cranmer of Satanism. No pedestrian word for word parody like that of the Luciferians, Crowley's Mass is both dramatic and original, a clear expression of the diabolic lex credendi. Prepared to effect what it signifies in the souls of its celebrants, it is a novus ordo in the real sense of the word, proclaiming the New Law in the first words of its Introiti ``Do what thou wilt shall be the whole of the Law!''

Crowley said, ``It replaces the moral and religious sanctions of the past, which have everywhere broken down, by a principle valid for each man and woman in the world, and self-evidently feasible.''

He also said, ``I resolved that my Ritual should celebrate the sublimity of the operation of universal forces without introducing disputable metaphysical theories. I would neither make nor imply any statement about nature which would not be endorsed by the most materialistic man of science. On the surface this may sound difficult; but in practice I found it perfectly simple to combine the most rigidly rational conceptions of phenomena with the most exalted and enthusiastic celebration of their sublimity.'' Thus does Satan win to himself the ``matter-of-fact'' man of science by classic gnostic dissembling!

That super-showman Anton Szandor LaVey, High Priest of the Church of Satan, in San Francisco, who proclaimed the opening of the Satanic Age on Walpurgisnacht, April 30, 1966, uttered a most disturbing comment in this connection. He said that pseudo-Satanists have always appeared more or less openly throughout history, but ``the real Satanist is not quite so easily recognized as such.'' He should know, and unfortunately the history of Gnosticism bears him out.

We have already noted its use of ambiguity and equivocation as means of propagating its errors among the orthodox faithful. Rhodes says, ``To invoke the essential demons, it was not even necessary to depart from the orthodox ritual. St. Cyprian and St. Ambrose, renowned for their power to put devils to flight, could be invoked, but in the contrary sense so that Satan was invited to approach instead of to depart. There may have been, as in other cases, practical as well as theoretical reasons for this technique. The priest could say his impious Mass in the presence of an orthodox congregation without their being any the wiser, but its potency would not be impaired thereby.

A whole group of practices grew up around this furtive Satanism. The Abbé Davot, drunken rector of La Voisin's parish church, was not alone in concealing figurines and incantations under the corporal on the altar when he said Mass. There were a very large number of his fellows who, although strongly suspected of Satanism, and worst escaped the fine meshes of the law because no positive evidence of any kind existed against them.''

Close to our own day, the Theosophist Bishop Leadbeater of the Old Catholic Church of Utrecht, Holland, was well-known for his attempts to tailor Catholic liturgy to gnostic specifications, adding, deleting, altering and perverting its ancient prayers into magical incantations designed to release the liturgy's ``magnetic force'' and accelerate human evolution. Harnessing what he termed ``the occult power of the Mass,'' Bishop Leadbeater had his ``revised liturgy'' ready for Easter Sunday, 1917. The Theosophist Annie Besant soon predicted this ``little-known movement called the Old Catholic, with the ancient ritual, with unchallenged Orders, yet holding itself aloof from Papal Obedience ... is a living Christian Church which will grow and multiply as the years go on, and which has a great future before it, small as it yet is. It is likely to become the future Church of Christendom `when He comes.'\,'' Such is the power of a changing liturgy.

Its possibilities could not have been absent from the mind of St. Pius V, who issued over 200 Bulls implementing the decrees of the Council of Trent at the very time the black Mass was cresting. It's significant that one of the first specifically revoked the permissions granted in some cases for the celebration of evening Masses as ``departing from the ancient customs of the Roman Church.'' ``Therefore desiring to purge the Church of this abuse,'' the bull threatened priests continuing the practice with automatic suspension. Yet another bull forbade the clergy to celebrate in any rite but their own, under the same penalty.

Closely following the canons of Trent, which had decreed that all Missals be purged of any ``superstitious or apocryphal prayers'' with texts rendered uniform and rubrics fixed ``according to the ancient usage of the Holy Roman Church'' so that ``the people be not scandalized by new or differing rites,'' St. Pius issued the Great Bull Quo Primum in 1570, at the summit of his sanctity and his pontificate. Canonizing once and for all the inviolable Mass of Tradition in the Latin rite, this document reads, ``We order and enjoin that this Our recently published Missal must never be added to, that none of it be omitted, and that none of it be changed under penalty of Our indignation.'' Prelates of whatever rank, furthermore, must not ``in celebrating Mass presume to introduce any ceremonies or recite any prayers other than those contained in this Missal,'' so delicate is the union of liturgy and doctrine.

This Mass, called Tridentine, is the work of the Holy Ghost, the very Spirit of Truth. The great Benedictine Abbot of Solesme Dom Gueranger said such liturgy ``is so excellent a thing, one must go to God himself to discover its principle.'' The precision of its language is ordered to defining clearly and absolutely, again and again, the exact identity of the God who is being adored, together with the exact nature of the Sacrifice this God is actually offering to Himself by the human agency of His ordained priest.

It is a standing monument to the commandment given to Moses: ``Thou shalt fear the Lord thy God, and shall serve him only'' -- the very command with which our Lord confronted Satan's temptation to worship him in the desert. Tu solus sanctus, Tu solus Dominus, Tu solus altissimus, Jesu Christe, in gloria Dei Patris!

There is no other. Simon Magus would have had difficulty with such words, but how glibly would he mouth some of the ambiguous passages of the Mass of Paul VI! For him, who would be the ``Lord God of creation,'' the Deus universi invoked at the offertory of the Mass of the Novus Ordo Seclorum? This Lord is not called the Creator but only the Lord of what has been created, of the created gifts about to be offered. Is he not the prince of this world whom our Lord told us was ``coming'' and of whom He said, ``In me hath not any thing?'' (Matt. l4:30).

This new service was in fact the work of one Annibale Bugnini, a prelate now known to be a Masonic adept. Whose Lord is his? It was our Lord who first asked His enemies, ``What think you of Christ? Whose son is he?'' To which they replied that he was the son of a man, a Jew ``David's,'' a lord of creation (Matt. 22:42). On the other hand, when our Lord asked His disciples, ``But whom do you say I am?'' they answered for all the faithful through the mouth of St. Peter that He was ``the Son of the living God'' (Matt. l6:15-16). The sword of division ever falls between these two irreconcilable answers of the sons of God and the sons of Belial. It falls today between two Masses.

Like his modern followers, Simon Magus would have welcomed the elimination of St. John's Gospel, which St. Pius V stipulated must be read at the close of all low Masses in perpetuity. Under direct inspiration of the Holy Ghost, John the Apostle must surely have had the gnostics of his day in mind when he composed such incomparable phrases as Omnia per ipsum facta sunt: et sine ipso factum est nihil quod factum est. ``All things were made through Him, and without Him was made nothing That has been made,'' closing on the impregnable ET VERBUM CARO FACTUM at which all Christendom once genuflected in unity of faith.

No one could attend Mass with supernatural sensibility without becoming aware of the satanic presence ever prowling about It, seeking to devour It. Even as he begins to say It at the foot of the altar, the celebrant complains to God, Quare me repulisti, et quare tristis incedo dum affllgit me inimicus? ``Why do I go about sad while the Enemy afflicts me?'' Judica me Deus, et discerne causam meam de gente non sancta! Ab homine inquo et doloso erue me! ``Give judgment for me, 0 God, and decide my cause against an unholy people. From the unjust and deceitful man deliver me!'' he pleads in the words of the Psalmist.

All the Church's liturgy is exorcism in the sense that it saves us from Satan's power through the Incarnation and death of the only-begotten Son of God. Before receiving Him in Holy Communion at Mass we ask in the Lord's Prayer, Libera nos a malo! In English this is translated, ``Deliver us from evil,'' but in the Latin, and especially in the Greek original, the words malo or tou ponerou can mean equally ``the Evil One.'' Against him, Adjutorium nostrum in nomine Domini qui fecit caelum et terram, says the ancient Mass unequivocally. Our help is in the name of the Lord who made not only heaven, but earth as well, material creation as well as the spiritual. Let gnostics apply elsewhere if they choose. ``The eternal Word who existed in the beginning,'' says the great liturgist Dom Guéranger, ``was made flesh, in time, and He lived among us in order to establish religion on the true worship whose visible symbols contain grace even as they signify it.'' During His life on earth this very Word openly accused the divinely instituted Jewish priesthood of departure from true worship, of ``transgressing the commandment of God'' in favor of their own man-made ``tradition'', the Kabbalistic precepts of the Babylonian Talmud.

The Book of Daniel gave witness that ``Iniquity came out from Babylon from the ancient judges, that seemed to govern the people'' (Dan. 13:5). And our Lord applied to them the words of Isalas: ``In vain do they worship Me, teaching doctrines and commandments of men''. (Matt. 15: 3,9).

In the prayer Supra quae, which unhappily forms no part of the new twentieth century Canons, the traditional Mass placed itself squarely in the line of worship which began with the holy Abel, begging God to hold it ``acceptable, as Thou didst deign to accept the offerings of Abel Thy just servant.'' Thus does it distinguish itself categorically from Cain's self-propelled offering of ``the fruits of the earth,'' to which God ``had no respect.''

Clearly, not all religions are good, nor may anyone worship as he pleases. Herein lies the essential evil of the Mass we have come to call ``black,'' now associated with black cats, obscenity and imprecations at midnight. Sacrilegious as they may be, these sensational aspects are merely peripheral. What makes a Mass ``black'' is, in the final analysis, its unacceptability to God. And it follows that what is not acceptable to God is readily so to Satan.

A black Mass might well be ``valid'' in the narrow sense of yielding the Sacrament. Certainly the ones Catherine de Medici ordered were considered so, nor does there seem to have been much doubt about those Fr. Davot offered so frequently before his unwitting parishioners during the reign of Louis XIV.

It has long been the common opinion of theologians that a duly ordained priest with the proper intention could say Cranmer's or even Luther's Mass validly, although these would certainly be illicit and sacrilegious. Where there is a proper Consecration the Church has never explicitly ruled otherwise. Years ago this was the basis of alarm on the part of some Catholics who became aware that certain Anglican ministers had got themselves ordained by schismatic bishops and were therefore probably capable of truly consecrating during Protestant services. It is also conceivable that one of the formulas for the double Consecration might be defective. This would confect the Sacrament validly under one species, but eliminate the Sacrifice by rendering impossible the mystical separation of Body and Blood which the double Consecration specifically effects.

In all these varieties of possibly valid ``black'' Masses, the intention of the celebrant is vital. Speaking in Florence in February, 1965, Archbishop Lefebvre was reported as saying, ``Certainly for a few years Luther said a valid Mass, while he was still not against the idea of Sacrifice, while he was more or less a Catholic; but later, when he rejected the Sacrifice, Priesthood and Real Presence, then his Mass was no longer valid. But how can a Mass be thus equivocal? It can't happen with the old Mass because the old Mass is so clear. The whole of the Offertory expresses what it is that we are doing. The Offertory is a definition of the Sacrifice of the Mass. That is why Luther was so opposed to the Offertory, because it was too clear, and for the same reason he made those changes in the Canon, so as to make it uncertain whether it was a reading or an action. And we know that the Consecration is a sacrificial action.'' The virtual elimination of the Offertory from the new ``Eucharistic Prayers'' of Vatican II speaks too loudly to require emphasis here.

If black Masses can be valid, they can also be licit, in the strictly juridical sense of conforming to ecclesiastical regulations or being permitted by established authority. They might even be licit without being valid, like those now permitted in many parts of Europe where they are said by the people without benefit of clergy. Inasmuch as the worst black Masses, in the sense of their utter unacceptability to God, need not ever be recognized as such by the unwary, it's all too true that they can, legalistically, satisfy the Sunday obligation!

From what the eye could detect, Cain's sacrifice must have seemed quite as reverent and holy as Adel's. It was certainly both ``legitimate'' and ``valid'', for as firstborn son no one had a better right than Cain to offer sacrifice under the natural dispensation then in effect. Indeed Scripture shows that he did in fact offer before Abel. Despite all liceity, however, God ``had no respect to Cain and his offerings.'' His was vain worship. Entirely ineffectual in bringing God's blessing down on the world, it actually provoked the world's first homicide and earned Cain perpetual exile. By their fruits are black Masses known, even as are their celebrants.

Over a hundred years ago Dom Guéranger wrote, ``If the sacrifice of the Mass were extinguished, we would not delay falling into the depraved condition in which peoples tainted with paganism found themselves, and such will be the work of the AntiChrist. He will seek every means of preventing the celebration of Holy Mass so that this great counter-weight may be overthrown and God will put an end to all things, having no longer any reason to keep them in existence. We can easily understand this, for since Protestantism, we notice far less strength in the heart of society. Civil wars have arisen bringing desolation in their wake, and that solely because the intensity of the sacrifice of the Mass is reduced. This is the beginning of what will happen when the devil and his followers will be unleashed over the world ... ''

Since Vatican II we note that just not the strength, but the very heart seems to have gone out of society. Placeat tibi, sancta Trinitas, obsequium servitutis meae, cries the ancient Mass to heaven for us, et praesta, ut sacrificium quod oculis tuae maiestatis indignus obtuli, tibi sit acceptabile! ``May the tribute of my worship be pleasing to You, most Holy Trinity!'' Grant that the ``Sacrifice I have offered may be acceptable!''

Source: Is the Black Mass Valid\footnote{\url{https://www.remnantnewspaper.com/Archives/archive-2006-black-Mass.htm}}

